\section{Introduction}\label{sec:intro}

\subsection{Goal}
The assignment asks for a self-hosted collaboration tool for a closed network: users sign in, see who
else is online, chat in shared channels or privately, and exchange files, without any external
service \cite{brief1,briefStreaprous}. Two constraints shape the whole design. First, only raw sockets
and the Python standard library may be used on the server side --- no Flask, Django or WebSocket ---
and the browser may run JavaScript only on the client side. Second, files must travel over
\emph{UDP}, so reliability, ordering, flow and congestion control have to be built at the
application layer.

\subsection{What was built}
\emph{Peer Chat} is a hybrid application with three kinds of processes:

\begin{itemize}
  \item a \textbf{tracker} (client-server part) that checks passwords, keeps the list of live peers
        with their user, IP address and ports, and can relay a frame when two peers cannot reach each
        other;
  \item one \textbf{peer daemon} per user (peer-to-peer part) that serves the web interface, sends
        chat frames \emph{directly} to the other peers over TCP and transfers files over UDP;
  \item an optional \textbf{reverse proxy} with virtual hosts and load balancing, inherited from the
        StreApRous skeleton and completed.
\end{itemize}

On top of the course skeleton the project adds an application protocol for chat over TCP, a reliable
file-transfer protocol over UDP (sequence numbers, selective and cumulative acknowledgements, a sliding
window driven by slow start and CUBIC, fast retransmit and an RFC~6298 retransmission timer),
AES-GCM encryption of every payload, a browser interface with unread counters and transfer progress
bars, and \FactTests{} automated tests, many of which start real processes and drop packets on purpose.

\subsection{Starting point and team work}
The code base started from the version of one team member (\todo{name}), which already had the UDP
file transfer with a simplified CUBIC window and AES-GCM encryption. A review against both assignment
documents found gaps --- no tests, no report, broadcast routed through the tracker, a tracker that
did not know the users, an unfinished proxy --- and several defects (\cref{sec:challenges}). This
report describes the system after those gaps were closed. \todo{state who did what; keep it
consistent with the table on the title page}

\subsection{Structure of this report}
\Cref{sec:requirements} maps the requirements to the implementation. \Cref{sec:architecture}
presents the architecture, \cref{sec:protocol} specifies the protocols and \cref{sec:lifecycles}
follows a login, a message and a file through the system. \Cref{sec:reliability} explains the
reliability and congestion-control mechanisms, \cref{sec:concurrency} concurrency and failure
handling, and \cref{sec:security} security. \Cref{sec:ui,sec:testing} cover the interface and the
evaluation, \cref{sec:challenges} the problems we met and how we solved them, and
\cref{sec:conclusion} the limitations.
